\paragraph{A diagnostic for distant evidence use.}
On the paired task, write the information-set error as
$E_{\mathrm{info}}=E_{\mathrm{first}}+E_{\mathrm{second}}$.
The first term measures four fair-marginal prediction errors; the second is
the sum of four deterministic conditional cross-entropies, averaged over
the 16 target-distributed first-round histories. Thus the stronger sufficient
condition for beating every one-call product predictor is
$E_{\mathrm{second}}/4<\log2-E_{\mathrm{first}}/4$.
This gives a specific learned quantity to improve through context adaptation.

A complementary local intervention flips one public parity constraint while
holding the visible information-set values fixed. Exactly one remaining
target flips. Orient its two logits so the required targets are respectively
zero and one, and denote them by $\ell_0,\ell_1$. With
$\Delta=\ell_1-\ell_0$ and $\operatorname{sp}(x)=\log(1+e^x)$,
\begin{equation}
 L_{\mathrm{flip}}
 =\tfrac12\{\operatorname{sp}(\ell_0)+\operatorname{sp}(-\ell_1)\}
 \ge \operatorname{sp}(-\Delta/2).
 \label{eq:evidence-response}
\end{equation}
Convexity proves the bound, with equality when $\ell_0+\ell_1=0$.
Small contrast imposes an error floor; a large common logit offset can still
produce high error despite a large contrast. We therefore distinguish response
strength from the nonnegative gap in Eq.~\eqref{eq:evidence-response}, and
check changes on unaffected coordinates. A fixed-history response probe is a
local diagnostic; the exact endpoint audit supplies the full history-averaged
error. Neither this elementary bound nor a response to one constraint alone
certifies broad long-context competence.
